\section{Executed experiments}
\label{sec:experiments}

\subsection{Eighteen runs, reported separately}

This section reports eighteen independent prototype runs in two rounds: the
design round \src{p1}--\src{p9}, covered first, and the extension round
\src{e1}--\src{e9}, covered in Section~\ref{sec:extension-runs}. All
eighteen are named after the source proposals listed in
Appendix~\ref{app:provenance}. They are reported separately and they are
\emph{not} pooled --- neither within a round nor across the two --- for a
reason stronger than caution.

All nine of the design-round runs ran on the same day in the same kind of
environment --- Python 3.13.5
on Linux x86-64, NumPy 2.3.5, SciPy 1.17.0, SymPy 1.14.0, and where pytest was
used, pytest 9.0.2 --- and five of them recorded the same seed value,
\code{20260914}. That coincidence makes the runs look comparable. They are not.
The same seed drives entirely different generators over entirely different case
sets: one run generates 30 cone and 60 quadratic cases, another 1{,}080 witness
tables and 300 CDCL(T) instances, another 100 cone and 49 invariant cases,
another 1{,}500 differential SAT cases, another 26 curated certificates. Three
runs (\src{p1}, \src{p8}, \src{p9}) record no seed at all.

\begin{remark}[The non-pooling rule]
The counts below are not commensurable. ``Cases'', ``checks'', ``assertions'',
``tests'', ``certificates'' and ``bundles'' are five different units, and several
runs use more than one. Summing them would produce a meaningless total, and any
such total would describe the prototypes' own constructed distributions rather
than any theorem corpus. Every number in this section carries its run label for
that reason.
\end{remark}

Two properties do hold uniformly, and they are the ones worth carrying forward.
First, no coefficient anywhere in the certificate corpus is stored as a
floating-point number: rationals are exact throughout, serialised as strings
such as \code{"125/18"} or as integer numerator/denominator pairs. A scan of
every certificate payload in all nine runs found exactly one float, and it is a
timing field. Second, every run's stored certificates replay under a
standard-library-only checker, verified with site packages disabled.

\subsection{The runs}

\subsubsection*{\src{p1} --- six certificate families, one uniform protocol}

275 pytest instances passed. A deterministic benchmark of 198 constructed
instances produced certificates in every case, and a separate
standard-library-only decoder accepted 198/198. Each case was searched three
times after numerical warm-up; the table reports medians of per-case medians.

\begin{table}[htbp]
\centering\small
\begin{tabular}{@{}lrrrr@{}}
\toprule
Component & Certified & Ablation & Search (ms) & Recheck (ms) \\
\midrule
Quadratic SOS            & 60/60 & 1 & 0.864 & 0.701 \\
Finite cone LP           & 17/17 & 0 & 1.699 & 0.295 \\
Bernstein box            & 18/18 & 1 & 0.685 & 0.389 \\
Polynomial recurrence    & 33/33 & --- & 2.513 & 0.403 \\
Integral affine witness  & 40/40 & --- & 0.074 & 0.027 \\
Univariate with zeros    & 30/30 & 0 & 1.930 & 0.505 \\
\bottomrule
\end{tabular}
\caption{\src{p1}. Ablations are restricted versions of these same prototype
mechanisms, \textbf{not} \grind{} and not any Lean tactic. A dash means no
ablation was run. Search time includes checking performed internally by the
search routine, so the recheck column is not the amount of checking hidden
inside the search column.}
\label{tab:p1}
\end{table}

The ablations: the quadratic one recognises only expanded nonnegative
combinations of even monomials; the cone one removes richer square candidates
or, for the domain example, the hypothesis products; the Bernstein one disallows
subdivision; the univariate one runs dyadic Bernstein search to depth 12 without
factorisation. The last row is the sharpest result in this run --- 30/30 against
0/30 --- and it is the experimental face of Theorem~\ref{thm:zero}.

Mutation tests alter targets, weights, hypotheses, recurrence base terms, affine
offsets, Bernstein split points, child nodes, claimed bounds, and Sturm chains,
and require rejection. Independent symbolic tests reconstruct Bernstein
expansions, expand quadratic certificates, and compare univariate root counts.
Limitations are preserved as regression cases: a degree-limited recurrence
search failing on a valid higher-degree answer, integer affine synthesis
rejecting a merely rational witness, Bernstein-only search failing at a
non-dyadic interior zero.

\subsubsection*{\src{p2} --- five families, per-family seeds}

31 unittest methods passed; several contain multiple seeded trials. 215
certificates were produced and all 215 replayed from storage.

\begin{table}[htbp]
\centering\small
\begin{tabular}{@{}lrrr@{}}
\toprule
Family & Cases & Exact replays & Time (s) \\
\midrule
List induction and generalisation & 6 & 6 & 0.005 \\
Polynomial additive recurrences & 65 & 65 & 0.241 \\
Affine witnesses with Farkas certificates & 63 & 63 & 0.352 \\
Finite-dictionary square certificates & 40 & 40 & 0.648 \\
Bernstein box certificates & 41 & 41 & 0.067 \\
\midrule
Total & 215 & 215 & \\
\bottomrule
\end{tabular}
\caption{\src{p2}. Seeds are per family: 420 recurrences, 421 witnesses,
422 squares, 423 Bernstein. Times are one local run and are not
hardware-normalised.}
\label{tab:p2}
\end{table}

The recurrence family comprises five named power sums and 60 seeded random
polynomial increments; it required 312 candidates in total, and the exact step
checks generated 247 counterexamples before the 65 certificates were accepted.
The 41 Bernstein problems used 82 leaves in total at maximum split depth four;
root-only checking certifies 25 of them and subdivision certifies all 41. The
six list proof traces contain 6, 3, 8, 10, 9 and 5 rewrites.

The run's own caveat is the important one: \emph{the 40 square-certificate
problems are constructed from the same finite dictionary the search uses}, so
their 100\% success rate tests implementation and rational reconstruction on a
designed distribution and must not be read as nonlinear proving coverage.
Recorded negative outcomes include the Motzkin polynomial finding no certificate
in the plain-square dictionary, the constant witness template failing on the
unbounded strip, and $(x-y)^2$ remaining \unknown{} at depth four.

\subsubsection*{\src{p3} --- assertion-counted mechanism ablations}

1{,}680 assertions passed; 24 certificate bundles were saved and all 24 replayed
under a standard-library-only checker. Seed \code{20260914}; total 2.27 s.

\begin{table}[htbp]
\centering\small
\begin{tabular}{@{}lrr@{}}
\toprule
Mechanism & Restricted & Enabled \\
\midrule
Cone certificates, curated goals & 0/10 diagonal-only & 9/10 full dictionary \\
Cone certificates, constructed & --- & 30/30 \\
Exact quadratic, constructed & --- & 60/60 \\
Exact quadratic, curated & --- & 4/4 \\
Bernstein box, 10 true inequalities & 3/10 unsplit & 8/10 adaptive \\
Ground Horn, 12 goals & 3/12 unsliced & 12/12 demand-sliced \\
List induction, 8 goals & 6/8 fixed accumulator & 8/8 generalised \\
\bottomrule
\end{tabular}
\caption{\src{p3}. Every ``restricted'' column is a restricted version of this
run's own implementation. None is \grind{}.}
\label{tab:p3}
\end{table}

Four negative controls were rejected by the quadratic path, seven corruptions by
the Bernstein checker, and eighteen by the induction checker. 250 random Horn
graphs and 500 comparisons confirmed that the sliced and unsliced modes agree
with a naive least-fixed-point computation. Both Horn modes still pay the full
index-construction cost, and the examples deliberately place decoys first.

The 60 generated quadratic examples are constructed from positive square
combinations and therefore carry an explicit selection bias; the ten cone goals
are deliberately selected mechanism tests, not an unbiased corpus.

\subsubsection*{\src{p4} --- the largest artefact corpus}

64 pytest tests passed. The deterministic experiment produced 1{,}494 rows and
1{,}238 saved successful artefacts, all replayed by checking modules that import
no numerical or symbolic search library. Seed \code{20260914}. This run is also
the only one that re-verified its corpus from an extracted archive copy and
confirmed that every family and status count matched the recorded run.

\begin{table}[htbp]
\centering\small
\begin{tabular}{@{}lr@{}}
\toprule
Family & Observed result \\
\midrule
Residue witness parameters & 1{,}080 checked: 831 proved, 249 impossible by gcd divisibility \\
Random CNF + integer difference logic & 300 oracle agreements: 189 UNSAT, 111 SAT \\
Sparse cone / equality-ideal certificates & 47 of 49; two deliberate \unknown{} \\
Exact rational nonnegative quadratics & 40/40 manufactured examples \\
Bernstein covering certificates & 12 of 13; the non-dyadic zero case \unknown{} \\
Pigeonhole CNF & 5 replayed UNSAT proofs \\
Structural equations & 3 synthesised or proved lemmas, 2 rejected false candidates \\
Induction ablation & 4 rows: 1 proved, 3 \unknown{} \\
\bottomrule
\end{tabular}
\caption{\src{p4}. These families are heterogeneous and must not be combined
into a general accuracy score.}
\label{tab:p4}
\end{table}

The structural ablation is this run's sharpest finding, and it is a genuine
$2\times2$ design rather than a single comparison:

\begin{table}[htbp]
\centering\small
\begin{tabular}{@{}ccc@{}}
\toprule
Proved append lemmas & Generalise accumulator & Outcome \\
\midrule
No  & No  & \unknown{} \\
No  & Yes & \unknown{} \\
Yes & No  & \unknown{} \\
Yes & Yes & \textsc{proved} \\
\bottomrule
\end{tabular}
\caption{\src{p4}. Only generalisation \emph{and} proved auxiliary lemmas
together close the accumulator reversal. The same goal and the same rewriting
engine are used in all four configurations, which is stronger evidence of
mechanism than a single all-features-enabled demonstration, and it shows why a
blanket ``try induction'' addition would miss part of the design.}
\label{tab:p4-ablation}
\end{table}

Two further details behind the aggregate counts. The 1{,}080 residue cases sweep
all $1\le m\le20$, $-3\le a\le5$ and $0\le b\le5$; the 249 impossible ones fail
the exact gcd divisibility condition. And for
$p(x)=(x-1/3)^2+10^{-4}$ on $[0,1]$ --- the interior-zero example given a
positive margin --- the prototype returns a tree with seven leaves, thirteen
nodes and maximum depth six, whereas $(x-1/3)^2$ itself reaches the depth limit,
exactly as Theorem~\ref{thm:zero} requires.

The run's own caveats are worth preserving: 64 is a test-runner count, not a
count of mathematical instances, since parametrised tests contain inner-loop
cases; forty of the 49 cone cases are manufactured from the search's own
certificate basis; the recorded timing is one in-process pass rather than a
repeated performance study; and the field named \code{certificate\_bytes}
includes the accompanying statement, not a proof payload alone. The run states
the non-pooling rule for itself in the sharpest available form: \emph{do not
divide 1{,}238 by 1{,}494 and advertise an overall accuracy}, because the
denominator mixes proof searches, model searches, a lemma batch, deliberately
false formulas, and overlapping ablation runs.

The pigeonhole family is recorded per instance, including the datum that is
least flattering to the method:

\begin{table}[htbp]
\centering\small
\begin{tabular}{@{}rrrrr@{}}
\toprule
Pigeons / holes & Decisions & Conflicts & Resolution nodes & Nonchronological jumps \\
\midrule
2 / 1 &  0 &  1 &   2 & 0 \\
3 / 2 &  1 &  2 &  10 & 0 \\
4 / 3 &  8 &  9 &  45 & 0 \\
5 / 4 & 30 & 31 & 160 & 0 \\
6 / 5 & 89 & 89 & 513 & 1 \\
\bottomrule
\end{tabular}
\caption{\src{p4}. The small number of non-chronological jumps is reported
rather than hidden. It is not evidence that this branch heuristic is competitive
on industrial SAT instances. Note also that this encoding --- at least one hole
per pigeon and at most one pigeon per hole --- differs from \src{p6}'s, so
Tables~\ref{tab:p6-dpll} and this one are not comparable.}
\label{tab:p4-pigeonhole}
\end{table}

\subsubsection*{\src{p5} --- expected-outcome benchmarking with negative controls}

73 pytest tests passed in 1.11 s. Seed \code{20260914}.

\begin{table}[htbp]
\centering\small
\begin{tabular}{@{}lll@{}}
\toprule
Experiment & Outcome & Interpretation \\
\midrule
Named nonlinear cases & 9 accepted / 12 &
Two true cases unresolved; one deliberately false case unresolved \\
Generated product-cone cases & 100 / 100 &
Constructed inside the covered cone; not representative sampling \\
False-polynomial controls & 0 accepted / 24 &
All \unknown{}; exact negative witnesses checked separately \\
Bernstein cases & 43 / 45 &
One true non-dyadic-zero case and one false case unresolved \\
Additive recurrence invariants & 49 / 49 &
Base and universal step identities checked \\
Horn growth family & 5 depths &
Both strategies prove all goals; search work differs sharply \\
Actual Lean comparison & Not run & No measured superiority over \grind{} \\
\bottomrule
\end{tabular}
\caption{\src{p5}.}
\label{tab:p5}
\end{table}

The run states its own most important qualification: the \textbf{109} accepted
nonlinear certificates are the nine named successes plus the 100 generated
cases, and they are \emph{not} 109 previously unsolved Lean theorems. The 24
negative controls are of the form $(x-a)^2-\varepsilon$ with rational
$\varepsilon>0$; the benchmark separately checks the exact negative value at $a$,
while the cone routine itself returns \unknown{} rather than a certified
refutation. Median recorded cone discovery time was about 21.5\,ms with a
maximum of about 103.6\,ms; median invariant synthesis about 4.1\,ms, including
interpolation and validation but excluding any Lean work. 58 exported
certificates --- nine named cone plus all 49 invariants --- were reloaded and
re-accepted by a separate validator.

\subsubsection*{\src{p6} --- proof-logging CDCL(T)}

This run's core is a differential test against an exhaustive reference.

\begin{table}[htbp]
\centering\small
\begin{tabular}{@{}lr@{}}
\toprule
Check & Result \\
\midrule
Random pure CNF cases & 1{,}000 \\
Random CNF + integer-difference cases & 500 \\
Satisfiable results & 1{,}278 \\
Unsatisfiable results & 222 \\
Differential mismatches & 0 \\
UNSAT certificates independently replayed & 222 \\
Truncated certificates rejected & 222 \\
Additional structural edge cases & 5 \\
Deliberately mutated theory certificates rejected & 2 \\
\bottomrule
\end{tabular}
\caption{\src{p6}. Formulas have three to eight Boolean variables, clauses of
three distinct variables with independent signs, and one to six times the
variable count in clauses; difference constraints use four integer variables and
constants in $[-3,3]$. The reference enumerates all Boolean assignments and uses
an all-pairs shortest-path consistency check; the solver under test uses CDCL
with negative-cycle theory conflicts. These are small exhaustive tests, not an
industrial benchmark.}
\label{tab:p6-sat}
\end{table}

The five edge cases are an empty conjunction, an empty clause, a tautology,
contradictory units, and repeated literals; the integer-complement boundary has
its own regression assertion.

\begin{table}[htbp]
\centering\small
\begin{tabular}{@{}lrrrr@{}}
\toprule
Instance & Variables & Clauses & DPLL decisions & CDCL decisions \\
\midrule
Padded core, 0 unused  &  2 &  4 &     1 &   1 \\
Padded core, 4 unused  &  6 &  4 &    31 &   5 \\
Padded core, 8 unused  & 10 &  4 &   511 &   9 \\
Padded core, 12 unused & 14 &  4 & 8{,}191 &  13 \\
Pigeonhole $4\to3$     & 12 & 22 &     8 &   8 \\
Pigeonhole $5\to4$     & 20 & 45 &    51 &  39 \\
Pigeonhole $6\to5$     & 30 & 81 &   374 & 148 \\
\bottomrule
\end{tabular}
\caption{\src{p6}. The baseline is \emph{this run's own} chronological
unit-propagating DPLL with the same static variable order and phase. It is not
\grind{} and not a production SAT solver.}
\label{tab:p6-dpll}
\end{table}

The run attaches the caveats itself, and they should be kept. The padded-core
family is an artificial diagnostic: it places unused Boolean variables before a
two-variable inconsistent core, which a production preprocessor would delete.
Its decision-count gap illustrates non-chronological backjumping; it is not
credible evidence of practical speedup. More pointedly, \emph{runtime gives a
soberer picture than decision counts}: on the largest pigeonhole case both
solvers took about 26\,ms despite their different decision counts, with replay
adding about 21\,ms. Operation counts are treated as diagnostic and no general
speed claim is made.

Ten polynomial search cases yielded seven certificates and three deliberate
\unknown{}s --- a false target, Motzkin at degree cap 6, and $(2x+y)^2$ at degree
cap 2. A further 100 randomly planted rational-square certificates exercised the
checker and were cross-checked symbolically; those are \emph{checker tests, not
100 additional search successes}, and the manually supplied discriminant
identity \eqref{eq:disc} is likewise separate from the seven discoveries. Six
invariant powers $p=0,\ldots,5$ yielded degrees 1 to 6 with final sample counts
3, 4, 5, 6, 7, 8, and each rejected the alterations $Q+1$ and $Q+n$, giving 12
mutation rejections on the synthesis side and 12 more under independent replay.

The saved corpus is 8 SAT/theory proof logs, 8 polynomial certificates (seven
discovered, one manually supplied), and 6 invariant certificates, all replayed
with site packages disabled. Three polynomial search files carry no certificate
and are explicitly skipped as \unknown{}. The replay program fails on a missing
corpus rather than reporting vacuous success.

\subsubsection*{\src{p7} --- component cases with corruption controls}

833 principal component test cases produced 891 accepted certificate checks;
330 deliberately corrupted certificates were rejected; 150 exact polynomial
cross-checks were run against a symbolic oracle; and 897 synthesised accumulator
formulas were compared against recursive execution. Seed \code{20260914}.

The run is explicit that \emph{cases and checks are different units}: some Horn
cases compare two successful derivations while others have no derivation, and
timing repetitions and archived demonstrations are excluded from the aggregate
check count. Five cone-search misses and two invariant-template misses are
expected outcomes, not wrong proofs.

26 representative certificates --- 9 cone identities, 5 lattice traces, 11
accumulator invariants, 1 Horn derivation --- all passed separate
standard-library-only replay. The lattice worker was additionally validated by
an exhaustive 539-case single-row grid checked against integer gcd, and by 120
random feasible systems whose free-basis samples were re-checked against every
row.

\subsubsection*{\src{p8} --- the cleanest dictionary ablation}

260 pytest tests passed, the highest count in the design round. No seed
recorded.

\begin{table}[htbp]
\centering\small
\begin{tabular}{@{}lrr@{}}
\toprule
Experiment & Restricted baseline & Extended mechanism \\
\midrule
Polynomial cones, 60 constructed instances & 17 certificates & 60 certificates \\
Positive rational-box instances, 21 & 0 unsplit & 21 subdivided \\
Affine accumulator templates, 121 & not measured & 121 synthesised and replayed \\
Out-of-template recurrences, 5 & not applicable & 5 \unknown{} outcomes \\
\bottomrule
\end{tabular}
\caption{\src{p8}. The cone row varies only the dictionary and holds the LP
backend fixed, which makes it the most cleanly isolated ablation in the
collection. It is still an ablation of this run's own implementation.}
\label{tab:p8}
\end{table}

The Horn workload is the largest: a useful chain of length 12 buried in 1{,}000
unrelated chains of length 24, each with an initial fact. Full closure fires all
24{,}012 rules; an indexed baseline that stops as soon as the goal is known still
fires 11{,}012 under the specified interleaved queue order; backward slicing
selects just the 12 useful rules.

The accompanying observation matters more than the ratio: \emph{proof
minimisation reduces even the full-closure derivation to the same 12 proof
steps}, so proof size alone would conceal the entire search-work difference.
Conversely the cold reverse index still scans all input rules and the demand
implementation still loads the initial fact set, so no sublinear end-to-end cold
performance is claimed. The run explicitly declines to infer a speedup factor,
noting a \emph{slower} early-goal run despite fewer firings.

Supporting evidence: 121 invariant cases each cross-checked against runtime
execution with four deterministic random-list checks per pair (484 concrete
recurrence checks); 30 seeded sparse-polynomial tests against a symbolic oracle;
40 randomised Horn programs comparing full and sliced success and replaying
their derivations; and monomial-enumeration counts checked against
$\binom{n+d}{d}$. Only two certificates were exported to files, the fewest of
the nine.

\subsubsection*{\src{p9} --- small, heavily mutation-tested}

69 pytest tests passed. 13 persisted bundles replayed with site initialisation
disabled: one ordered bank of five inductively proved lemmas with a
specialisation proof, four Horn proofs, six cone certificates, one box
certificate, and one witness with its two side-condition proofs. These are
bundles, not 13 independent Lean theorems. No seed recorded.

\begin{table}[htbp]
\centering\small
\begin{tabular}{@{}lrrrr@{}}
\toprule
Depth & Forward facts & Backward proof nodes & Forward cap & Forward status \\
\midrule
 4 &    16 &  5 & 5{,}000 & Certified \\
 8 &   256 &  9 & 5{,}000 & Certified \\
12 & 4{,}096 & 13 & 5{,}000 & Certified \\
16 & 5{,}000 & 17 & 5{,}000 & Cap reached \\
\bottomrule
\end{tabular}
\caption{\src{p9}. Neither column is an implementation of \grind{}. The forward
baseline stops as soon as the target is produced, which is why its counts are
$2^d$ rather than $2^{d+1}-1$.}
\label{tab:p9-horn}
\end{table}

Mechanism outcomes: accumulator-lemma discovery enumerates 41 candidates,
filters against 49 sample environments, and the 41st enumerated term
$\mathrm{app}(\mathrm{rev}(xs),a)$ is the one that receives an induction
certificate; the quartic nonnegativity case builds 280 candidate columns and 35
coefficient equations and returns a three-term certificate; the witness search
reaches candidate 23, $w(x)=x^2+1$, with three candidates reaching universal
certificate search; the Bernstein case $x^2-x+3/10$ on $[0,1]$ needs one split
and yields two leaves each with exact lower bound $1/20$.

The deliberate control is $(x+y+1)^2$: the restricted dictionary finds no
certificate, while a manually supplied one-square certificate is accepted. This
run therefore has a \emph{checker more general than its searcher}, which it uses
explicitly as a guard against an inflated ``SOS-complete'' claim.

This run has the most systematic mutation testing in the design round, with
parametrised
bad-input matrices for cone, box, Horn and induction certificates, including
rejection of a base-case use of the induction hypothesis, invalid generalisation
metadata, rigid-variable intrusion, and reference to an unavailable future
lemma.

\subsection{Four Horn experiments that must not be merged}

Four runs measured demand-directed versus forward Horn search. They used four
incompatible setups, and the temptation to combine their tables should be
resisted.

\begin{table}[htbp]
\centering\small
\begin{tabularx}{\textwidth}{@{}lX@{}}
\toprule
Run & Setup and headline observation \\
\midrule
\src{p5} & Depth sweep 4/8/10/12/14 on $P(a)$, $P(x)\to P(g(x))$,
$P(x)\to P(f(x))$; no distractors; no forward cap. Forward generates
$2^{d+1}-1$ facts (31, 511, 2{,}047, 8{,}191, 32{,}767) against $d+1$ demand
atoms; both proofs have $d+1$ nodes. Timings reported (846.283\,ms against
0.478\,ms at depth 14). \\
\src{p9} & Same rule family, depths 4/8/12/16, but the forward baseline
\emph{stops when the target appears}, giving $2^d$ facts; a 5{,}000-fact cap is
reached at depth 16. No timings reported. \\
\src{p3} & 12 curated goals with decoys placed first; 3/12 unsliced against
12/12 sliced; 250 random graphs and 500 agreement comparisons. \\
\src{p8} & Fixed useful chain of length 12 with a distractor sweep to 1{,}000
chains; 24{,}012 / 11{,}012 / 12 firings for full closure, early-goal stop, and
demand slicing. Timings explicitly disowned. \\
\bottomrule
\end{tabularx}
\caption{The same mechanism, four incomparable measurements. \src{p5} and
\src{p9} differ only in the forward baseline's stopping rule, and that alone
changes the reported growth from $2^{d+1}-1$ to $2^d$.}
\label{tab:horn-four}
\end{table}

That last point is the useful methodological lesson: a factor-of-two difference
in a headline growth curve came entirely from a baseline policy decision, not
from the mechanism under test.

\subsection{The extension round: nine more runs, still not pooled}
\label{sec:extension-runs}

A second round of nine independent proposals, labelled \src{e1} to \src{e9},
was prepared against the merged draft. Their runs are reported here separately
from the first nine and separately from each other. The non-pooling rule applies
with the same force --- and with an additional reason: the two rounds attack
overlapping problems, so a combined total would double-count the same
mathematics under different unit names.

The environmental coincidence repeats almost exactly. All nine ran on Python
3.13.5 on Linux 6.18.44 x86-64 with SymPy 1.14.0, and five of them
(\src{e1}, \src{e2}, \src{e3}, \src{e4}, \src{e6}) recorded the seed
\code{20260915} --- one day after the \code{20260914} shared by five of the
design-round runs. \src{e5} used \code{2026091507}, \src{e9} used
\code{681437}, and \src{e7} and \src{e8} record fixed seeds without a shared
value. As before, the shared seed drives entirely different generators over
entirely different case sets, and the coincidence carries no comparability.

\begin{table}[htbp]
\centering\small
\begin{tabularx}{\textwidth}{@{}llX@{}}
\toprule
Run & Headline unit & Recorded outcome \\
\midrule
\src{e1} & 124 runs / 62 problems &
74 proved, 34 refuted, 16 \unknown{}; 108 accepted certificates, all 108
replayed with site packages disabled and neither the search module nor SymPy
imported (0.052\,s); 40 unit tests. \\
\src{e2} & 65 accepted certificates &
Four families (25 two-sided, 24 invariant, 15 Gosper, 1 telescoper);
1{,}100 assertions; 188 mutations rejected; 527 pointwise binomial regressions
over $n=0,\ldots,30$; 100 differential substitution cases; 7 expected
\unknown{}s. \\
\src{e3} & 317 cases / 313 bundles &
178 proved, 135 refuted, 4 \unknown{}; 240 independent finite-state oracle
comparisons; 210 diagnostic moment evaluations; 25 targeted invalid mutations,
all rejected. \\
\src{e4} & 82 corpus entries / 73 certificates &
Machine lane 37 proved, 28 refuted, 1 \unknown{}; IPC lane 8 countermodels and
8 \unknown{}; 222 generated local Lean obligations, none compiled; largest
certificate 915 bytes; whole corpus 0.842\,s over three repeats; 151 pytest
items. \\
\src{e5} & 59 cases &
38 certified, 17 counterexamples, 4 \unknown{}, across 36 linear, 13 ideal and
10 binomial-sum cases; 818 mutations generated, 0 accepted; 117 pytest items. \\
\src{e6} & 82 attempts &
Ideal lane 27 certificates / 2 \unknown{}; weighted lane 23 closures and 25
distinguishing words; telescoping lane 3 certificates / 2 \unknown{}; 78
certificates replayed under a standard-library-only checker; 318 pytest
items. \\
\src{e7} & 135 replayed records &
113 constant-action invariant, 14 ideal-preservation, 4 ranking, 1
indexing-algebra, 3 orbit-counterexample records; 31 corruptions rejected; 29
test methods; whole run about 1.8\,s. \\
\src{e8} & 422 stored certificates &
51 constant-multiplier, 14 polynomial-multiplier, 136 finite-algebra, 221
integer-projection; all replay with the search entry points replaced by
exceptions; 3{,}400 parameter valuations checked against a complete concrete
oracle (1{,}368 feasible, 2{,}032 infeasible, no disagreement); 21 test methods
containing 19 explicit mutation cases. \\
\src{e9} & 245 accepted objects &
54 matrix closures, 115 separating words, 23 telescopers, 6
common-left-multiple identities, 47 singularity/seed plans; 245 specified
invalidating mutations, all rejected; 1{,}140 exhaustive word evaluations; 483
additional recurrence checks; 41 unit tests; 5.558\,s for the corpus and
0.339\,s for replay with SymPy absent from the replay process. \\
\bottomrule
\end{tabularx}
\caption{The extension round. Nine incommensurable ledgers: ``runs'',
``cases'', ``bundles'', ``certificates'', ``records'', ``attempts'' and
``objects'' are seven different units, and no row may be added to any other.}
\label{tab:extension-runs}
\end{table}

\subsubsection*{The ablations that carry the argument}

Three of these runs contain a controlled comparison that is more informative
than any of the totals, because in each case the \emph{same} problem set is run
through two lanes that differ in exactly one certificate language.

\src{e1} ran sixteen groups: eight problem families, each in a linear-span mode
and an ideal mode. The two modes agree everywhere except in two families. On
\emph{coupled nonlinear} and \emph{nonlinear equivalence}, eight cases each, the
linear lane returns \unknown{} on all sixteen and the ideal lane proves all
sixteen. That is the entire content of Section~\ref{sec:span-insufficient}
reduced to a measurement: no span budget closes those systems, and the failure
is not a fuel problem.

\src{e2} ran 24 invariant cases through the same feature set twice. The subspace
method certifies 24; restricting the certificate language to conserved
quantities certifies 3. The other 21 are true, and the restricted lane simply
cannot say so.

\src{e6} separates its weighted lane's outcomes rather than reporting a success
rate: of 48 attempts, 23 produce closure certificates and 25 produce
distinguishing words. Both are results. Counting only the first would report a
48\,\% success rate for a worker that actually answered every question it was
asked.

\subsubsection*{\src{e7}'s composition table, and how to read it}

One table deserves separate handling because it is the most quotable number in
the extension round and the most easily misquoted. Over four finite hole
grammars held fixed in size and order, a flat whole-term policy and a
contract-directed local policy both return the same term:

\begin{table}[htbp]
\centering\small
\begin{tabular}{@{}rrrrl@{}}
\toprule
Distractors & Flat complete terms & Local checks & Local complete terms & Hole sizes \\
\midrule
0  & 24     & 13 & 1 & 1, 2, 2, 8 \\
4  & 2{,}148  & 29 & 1 & 5, 6, 6, 12 \\
8  & 14{,}384 & 45 & 1 & 9, 10, 10, 16 \\
12 & 50{,}940 & 61 & 1 & 13, 14, 14, 20 \\
\bottomrule
\end{tabular}
\caption{Same grammar, same within-hole ordering, different composition
strategy. Local checks and complete terms are different work units and are not
comparable to each other.}
\label{tab:composition}
\end{table}

The distractor ordering is deliberately unfavourable and is held fixed across
both policies. The largest case was additionally checked on all lists of length
at most six over each of two two-element alphabets, at ordinary nearby indices
and at $\pm2^{80}$, agreeing with an independent non-fold reference on all
3{,}570 checks --- with the symbolic replay, not those observations, carrying
the general claim. This is a demonstration that a cross-product disappears under
factored contracts. It is not a measured $50\,940$-fold speedup, not robustness
under all candidate orders, and not evidence that the original Leant query has
been fixed.

\subsubsection*{What the re-run added}

Unlike the design round, whose authoring environments were unavailable here,
every extension suite could be executed in the merge environment. All nine were
re-run, and every recorded count was reproduced:

\begin{table}[htbp]
\centering\small
\begin{tabular}{@{}llll@{}}
\toprule
Run & Command & Recorded & Observed \\
\midrule
\src{e1} & \code{pytest -q test\_closure.py}    & 40 tests & 40 passed \\
\src{e2} & \code{pytest -q test\_regressions.py}& 13 methods & 13 passed, 69 subtests \\
\src{e3} & \code{pytest -q tests}               & 12 methods & 12 passed, 313 subtests \\
\src{e4} & \code{pytest -q tests}               & 151 tests & 151 passed \\
\src{e5} & \code{pytest -q tests}               & 117 items & 117 passed \\
\src{e6} & \code{pytest -q tests}               & 318 passed & 318 passed \\
\src{e7} & \code{pytest -q test\_extensions.py} & 29 methods & 29 passed, 140 subtests \\
\src{e8} & \code{python -S -m unittest test\_prototypes} & 21 methods & 21 passed \\
\src{e9} & \code{pytest -q tests}               & 41 tests & 41 passed \\
\bottomrule
\end{tabular}
\caption{Re-run in the merge environment --- Python 3.14.4, NumPy 2.4.4, SciPy
1.17.1, SymPy 1.14.0, pytest 9.0.3, Windows. Recorded in
\texttt{results/extension-suites-rerun.json}.}
\label{tab:extension-rerun}
\end{table}

Two observations, neither of them a defect. \src{e8} runs under \code{python -S},
which confirms its standard-library-only claim by making site packages
unreachable. \src{e9}'s suite must be run from \code{prototype/}; from the
package root it fails collection with
\code{ModuleNotFoundError: forge\_summaries}. That is a path convention.

The re-run is a reproduction on \emph{different} interpreter and library
versions than the authoring environment, not a bit-identical replay, and no
timing was compared. It establishes that the suites pass and that the recorded
counts are accurate. It does not establish that any checker is correct, and it
is not a Lean result.

\subsection{Lean status}

\subsubsection*{What the eighteen runs recorded}

None of the eighteen compiled any Lean source. Every one of them recorded the
same reason --- no \code{lean} or \code{lake} executable in the authoring
environment --- under several different filenames and in more than one format.
Consequently none performed a kernel check, an axiom audit, or any comparison
against \grind{}, Aesop, \code{nlinarith}, LeanHammer, or any other tactic.
Twenty-one \code{.lean} files across the design round, totalling about
3{,}000 lines, and twenty across the extension round, totalling about
3{,}100, were \emph{candidate source}. That the two rounds independently
produced almost exactly the same quantity of uncompiled Lean is itself worth
noticing: the bottleneck was never the amount of Lean anyone was willing to
write.

\subsubsection*{What this merge was able to add}

The merge environment has \code{elan}, and installing the pinned
\code{leanprover/lean4:v4.34.0} made a first, narrow compilation check possible.
Three of the design-round files in the merged tree import nothing beyond Lean
core and can therefore be elaborated without a built Mathlib:

\begin{table}[htbp]
\centering\small
\begin{tabular}{@{}llr@{}}
\toprule
File & Result & Elapsed (s) \\
\midrule
\src{p3} \code{lean/DesignAPI.lean}   & elaborated, no errors & 211.4 \\
\src{p6} \code{lean/Contracts.lean}   & elaborated, no errors & 5.3 \\
\src{p9} \code{lean/Structural.lean}  & elaborated, no errors & 5.1 \\
\bottomrule
\end{tabular}
\caption{Lean v4.34.0 (commit \texttt{293d5d0c}), recorded in
\texttt{results/lean-core-elaboration.json}.}
\label{tab:lean}
\end{table}

This confirms that the proposed runtime and scheduler contract, the
\code{CertificateSpec} soundness contract, and the Mathlib-free structural list
proofs are well-typed Lean 4.34.0 rather than merely plausible-looking source.

It confirms nothing else. It is not a transitive axiom audit. It says nothing
about the thirteen files in the merged tree that import Mathlib and need a
built project, or about the two that merely import their siblings. It does not establish that any certificate checker is correct.
And it is not a comparison against any Lean tactic --- successful elaboration of
a hand-written example proves that the example type-checks, not that automation
found it.

\subsubsection*{Ten more files, from the extension round}

The same check was then applied to the extension proposals. Of their twenty
source files, two are aggregators that import their siblings, eight import
Mathlib, and ten import nothing beyond Lean core. All ten elaborate:

\begin{table}[htbp]
\centering\small
\begin{tabular}{@{}llr@{}}
\toprule
File & Result & Elapsed (s) \\
\midrule
\src{e1} \code{lean/Core.lean}                          & elaborated & 82.1 \\
\src{e2} \code{lean/ForgeExtensions/Reachability.lean}  & elaborated & 175.8 \\
\src{e3} \code{lean/BridgeSpec.lean}                    & elaborated & 220.7 \\
\src{e4} \code{lean/Reachability.lean}                  & elaborated & 14.7 \\
\src{e5} \code{lean/ForgeExtension/Core.lean}           & elaborated & 9.4 \\
\src{e6} \code{lean/ProofPrinciples.lean}               & elaborated & 10.2 \\
\src{e7} \code{lean/Specimens.lean}                     & elaborated & 9.5 \\
\src{e8} \code{lean/CoreSoundness.lean}                 & elaborated & 8.8 \\
\src{e8} \code{lean/IndexingSketch.lean}                & elaborated & 30.8 \\
\src{e9} \code{lean/WordSimulation.lean}                & elaborated & 2.2 \\
\bottomrule
\end{tabular}
\caption{Lean v4.34.0 (commit \texttt{293d5d0c}), recorded in
\texttt{results/lean-extensions-elaboration.json}.}
\label{tab:lean-extensions}
\end{table}

Seven of the ten additionally print \code{\#print axioms} results, and every
printed line reads \emph{does not depend on any axioms} --- covering the
reachability and run invariants of \src{e2}, \src{e4}, \src{e6} and
\src{e8}, the tree-cover lemmas of \src{e8}, the fold-indexing and ranked-call
lemmas of \src{e7}, and the word-simulation lemmas of \src{e9}. This is a
stronger statement than the design-round result, where two of the three
elaborated theorems depend on \code{propext}.

The same three disclaimers apply, and one more besides. \src{e8}'s
\code{IndexingSketch.lean} elaborates, and its author had marked it
\emph{uncompiled} and \emph{not a substitute for the original query}. Both
statements are now true at once: the file type-checks, and it still proves an
equation between two ordinary-list definitions rather than anything about
Church-encoded inputs. Elaborating a specimen does not upgrade what the specimen
says.

\subsubsection*{Three more, written by the merge}

Those ten files overlap heavily. Between them they contain six spellings of one
reachability-invariant theorem and five of one word-fold preservation theorem,
differing in constructor names, argument order, and whether the initial state is
a predicate or a point. Merging them means proving each once and keeping what is
actually distinct, so the merged tree gains three files:

\begin{table}[htbp]
\centering\small
\begin{tabularx}{\textwidth}{@{}lXr@{}}
\toprule
File & Merged from & Elapsed (s) \\
\midrule
\code{Forge/Closure/Principles.lean} &
all nine, for the reachability and word-fold theorems; \src{e3} and \src{e4}
for the two-machine relation; \src{e3}, \src{e6} and \src{e9} for the
commuting step map; \src{e1} for the left fold; \src{e7} for well-founded
descent & 6.6 \\
\code{Forge/Closure/Covers.lean} & \src{e8}, the tree half of its soundness
specimen & 244.4 \\
\code{Forge/Closure/Indexing.lean} & \src{e7}'s Church half and \src{e8}'s
\code{List}/\code{Int} sketch, which proved the same equation under different
names & 3.8 \\
\bottomrule
\end{tabularx}
\caption{Recorded in
\texttt{results/lean-closure-elaboration.json}. All three elaborate, and all
twelve of their theorems print \emph{does not depend on any axioms}.}
\label{tab:lean-closure}
\end{table}

Sixteen files therefore elaborate in total: six in the merged tree and ten in
the extension proposals, with ten of the sixteen reporting no axiom
dependencies. The merged \code{Indexing.lean} inherits \src{e8}'s scope note
verbatim rather than quietly widening it.

\subsection{What the evidence establishes}

It establishes that the implementations actually synthesise certificates in the
stated fragments; that certificate replay is separable from the numerical and
symbolic search libraries; that the specified mutation and corruption tests are
rejected; that several mechanisms solve problems their own restricted precursors
cannot; that all nine extension suites re-run and reproduce their recorded
counts on a different interpreter and different library versions; and that
sixteen Lean artefacts --- three from the design round, ten from the extension
round, and three the merge wrote by deduplicating those ten --- type-check under
the pinned toolchain, ten of them printing \emph{does not depend on any axioms}
for every theorem they expose.

It does not establish that \forge{} is implemented as a Lean tactic, that the
Mathlib-dependent candidate scripts compile, that any certificate is checked by
Lean's kernel, that the methods outperform existing Lean tactics, or that any
success rate transfers to user developments. Those claims require different
experiments, described next.

Two structural limitations of the evidence deserve to be named rather than
buried. Manufactured examples favour the grammar that generated them --- most
visibly for \src{p2}'s 40 square certificates and \src{p5}'s 100 generated cone
cases, both of which are drawn from the very dictionary their search covers. And
search and replay share substrate: the same sparse polynomial and term
representations appear on both sides in every prototype. Search-free replay is a
meaningful independence property, but it is not an independently implemented
formal verification of that representation. The strongest partial answers to
this are the checkers that run algorithmically \emph{opposite} to their
producers --- expanding a claimed Bernstein basis rather than recomputing
coefficients, re-deriving a clause by independent unit propagation rather than
replaying a resolution chain, recomputing a Sturm chain rather than trusting a
stored root count, verifying $\det U=\pm1$ rather than calling the elimination
that produced $U$.

The extension round sharpens that last point into a three-level scale rather
than a yes/no property, and it is worth stating explicitly because the levels
are easy to conflate. At the strongest level the checker performs a genuinely
different computation from the search: \src{e7}'s ideal lane does not verify a
Gr\"obner basis, re-run Buchberger, or trust a zero-remainder flag, and
\src{e9}'s checkers multiply rational matrices and compare, where the searches
solved nullspaces. At an intermediate level the search entry points are
\emph{replaced by exceptions} during replay --- \src{e8}'s invariant and finite
lanes --- which proves separation but leaves input-table semantics shared. At
the weakest level, \src{e8}'s integer projection reconstructs the witness
program with the same assembler the search used, so replay catches a modified
receipt but not an arithmetic mistake common to both. The proposal says so
itself and mitigates it by a different route, an exhaustive concrete oracle over
a provably sufficient range. A reader comparing ``all certificates replay''
claims across the eighteen prototypes should ask which of these three things
the sentence means.

One more limitation is specific to the second round and does not appear in the
first. Several extension corpora are \emph{constructed positives}: \src{e7}'s
twelve power-curve fixtures have known invariants, \src{e9}'s forty comparison
models are invertible coordinate changes of random systems, and \src{e8}'s
forty-eight affine-diagonal instances are built to preserve the diagonal. These
exercise polynomial multiplication, substitution and coefficient handling. They
are not held-out benchmarks of discovery performance, and the proposals label
them as such rather than counting them as solved theorems. Likewise \src{e7}'s
seventy-two affine differential systems include forty-eight with a zero relation
space; those are useful controls, not additional nontrivial results.
